\begin{theorem}
Let $C>1$ be fixed.  Interpret the paper's integer $Cs$ as
$k_s=\lceil Cs\rceil$.  Then, for all sufficiently large $s$,
\[
 r(s,k_s)\geq\left(2^{1-1/(2C)}\right)^s.
\]
The ceiling changes only a bounded rounding error and is the precise integer
convention used here.
\end{theorem}
\begin{proof}
Fix $C>1$.  Since $k_s=\lceil Cs\rceil\geq Cs>s$, there is a threshold
$S_0$ such that for every integer $s\geq S_0$ we have $4\leq s\leq
k_s$.  Fix such an $s$ and abbreviate $k=k_s$.  By
\noderef{F2ForwardIndependentBound}, there is a loopless $T_s$-free digraph
$D=D_2(s)$ with
\[
N=2^{2s-3}-2^{s-2}
  =\bigl(1-2^{-(s-1)}\bigr)2^{2s-3}
\]
vertices and
\[
 \operatorname{fwi}_k(D)\leq
 \sum_{t=1}^{s-1}\binom kt
 2^{(s-1)(t+k)-\binom{t+1}{2}}.
\]

For $1\leq t\leq s-1$, complete the square in the exponent:
\[
(s-1)(t+k)-\binom{t+1}{2}
 = (s-1)k+\frac12(s-\tfrac32)^2
   -\frac12\bigl(t-(s-\tfrac32)\bigr)^2.
\]
Thus this exponent is at most $(s-1)k+\frac12(s-\tfrac32)^2$.
Since $\sum_{t=1}^{s-1}\binom kt\leq2^k$, it follows that
\[
 \operatorname{fwi}_k(D)\leq
 2^k2^{(s-1)k+\frac12(s-\frac32)^2}
 \leq 2^{sk+s^2/2}.
\]

Apply \noderef{RandomPermutationReduction} to $D$.  It gives a loopless
$K_s$-free graph $G$ on the same $N$ labelled vertices such that
\[
i_k(G)\leq\left(\frac ek\right)^k\operatorname{fwi}_k(D)
 \leq\left(\frac ek\right)^k2^{sk+s^2/2}.
\]
If $i_k(G)=0$, then $G$ has neither a clique of order $s$ nor an
independent set of order $k$.  Hence, by the defining property of
\noderef{RamseyNumber}, $r(s,k)>N$, and this is enough because $N$ is
eventually larger than $2^{(1-1/(2C))s}$.

Suppose now that $i_k(G)>0$.  As an integer count, $i_k(G)\geq1$; put
$p=i_k(G)^{-1/k}$.  Then $0<p\leq1$ and $p^ki_k(G)=1$, while $k\geq1$.
Therefore all hypotheses of \noderef{SamplingDeletion} hold: $G$ is
loopless and $K_s$-free by construction, $p$ lies in $[0,1]$, and the
sampling count inequality holds with equality.  That node gives
\[
r(s,k)>pN-1.
\]
Taking the inverse $k$-th root of the preceding upper bound on $i_k(G)$
gives
\[
p\geq \frac{k}{e}2^{-s-s^2/(2k)}.
\]
Consequently,
\[
pN\geq
\bigl(1-2^{-(s-1)}\bigr)\frac{k}{e}
 2^{s-3-s^2/(2k)}.
\]
Let $T=2^{(1-1/(2C))s}$.  Because $k=\lceil Cs\rceil\geq Cs$,
\[
s-3-\frac{s^2}{2k}
 =\left(1-\frac1{2C}\right)s-3
   +\frac{s(k-Cs)}{2Ck}
 \geq\left(1-\frac1{2C}\right)s-3.
\]
It follows that
\[
pN\geq
T\,\bigl(1-2^{-(s-1)}\bigr)\frac{k}{8e}.
\]
The factor multiplying $T$ tends to infinity, because $k\geq Cs$ and
$1-2^{-(s-1)}\to1$.  Increase $S_0$ so that for all $s\geq S_0$ this
factor is at least $2$ and $T\geq1$.  Then $pN-1\geq2T-1\geq T$, so the
sampling conclusion yields $r(s,k)>T$, and hence $r(s,k)\geq T$.
Since $k=k_s$, this is exactly the stated assertion for every $s\geq S_0$.
\end{proof}
